\documentclass[11pt,a4paper]{article}
\usepackage[T1]{fontenc}
\usepackage[utf8]{inputenc}
\usepackage[french]{babel}
\usepackage[margin=2.4cm]{geometry}
\usepackage{booktabs}
\usepackage{array}
\usepackage{enumitem}
\usepackage{hyperref}
\setlist{nosep}

\title{OpenMedReg-AI --- Première campagne de mesure\\
\large Deux juges locaux sur le cas fictif \texttt{fictif-therm-t1}}
\author{Rapport rédigé par un agent (CLAUDE), non relu par un humain}
\date{11 novembre 2026}

\begin{document}
\maketitle

\section*{Avertissement}

Ce rapport décrit une mesure, pas une validation. L'outil aide à la
relecture ; il ne délivre pas d'avis de conformité. Le cas mesuré est un
dossier \emph{fictif} rédigé par une IA, et sa vérité terrain (44
annotations) porte le statut \texttt{non\_relu} : aucun humain compétent en
affaires réglementaires ne l'a encore validée. Les chiffres ci-dessous
mesurent donc l'accord entre un modèle et une annotation produite par une
autre IA. Ils ne permettent ni de conclure à la fiabilité de l'outil, ni de
classer les deux modèles entre eux.

\section{Objet}

Première exécution réelle du moteur d'analyse d'écarts. Jusqu'ici, aucun
juge n'avait tourné ailleurs que contre des doublures de test. Deux modèles
ouverts, servis par Ollama sur un GPU loué, ont jugé le même dossier avec
les mêmes consignes. Le juge par API propriétaire, prévu comme point de
comparaison haut, n'a pas encore été mesuré.

\section{Protocole}

\begin{description}[leftmargin=3.2cm,style=nextline]
\item[Cas] \texttt{eval/cases/fictif-therm-t1} : un dossier technique
  fictif en français et 44 unités feuilles du périmètre \texttt{anx\_II}
  du corpus (texte consolidé du règlement, fichier
  \texttt{corpus/public/mdr/}, provenance dans \texttt{SOURCE.yaml}).
  Vérité terrain : 23 \texttt{conforme}, 6 \texttt{non\_conforme},
  7 \texttt{non\_trouve}, 8 \texttt{non\_applicable} ; 13 lacunes voulues.
\item[Empreintes] Début du SHA-256 de \texttt{dossier.md} :
  \texttt{85fe3e1deb403d8e}.\\
  Début du SHA-256 de \texttt{ground\_truth.yaml} :
  \texttt{ea721ec6b4ca358c}.
\item[Code] Commit \texttt{e2b1fc5} (moteur, consignes et juge inchangés
  depuis \texttt{ee6aec1}). Le second modèle a été lancé depuis une copie
  de ce commit, pour ne pas dépendre de modifications en cours dans le
  dossier de travail. Aucun masquage.
\item[Appels] Un appel par unité ; le dossier entier est fourni à chaque
  appel ; température 0 ; fenêtre de contexte de 16\,384 jetons ; sortie
  contrainte par un schéma JSON.
\item[Règle de citation] Un verdict dont la citation du règlement ou du
  dossier n'est pas retrouvée mot pour mot est ramené à
  \texttt{non\_trouve} avant tout calcul.
\item[Modèles] \texttt{qwen3.8:27b} (27,3 milliards de paramètres) et
  \texttt{gemma4:31b} (30,7 milliards), quantification Q4\_K\_M,
  Ollama 0.40.16. Mode de raisonnement laissé à sa valeur par défaut
  (actif) pour les deux. Détail dans \texttt{modeles.txt}.
\item[Matériel] Une tranche de carte graphique louée (24 Go annoncés,
  23,6 Gio vus par Ollama), modèles entièrement chargés sur la carte.
\end{description}

\section{Résultats}

Les intervalles sont des intervalles de Wilson à 95\,\%.

\begin{center}
\begin{tabular}{@{}lccc@{}}
\toprule
 & \texttt{qwen3.8:27b} & \texttt{gemma4:31b} & Repère « toujours \texttt{non\_trouve} » \\
\midrule
Exactitude & 33/44 (75\,\%) & 38/44 (86\,\%) & 7/44 (16\,\%) \\
Intervalle & 61 à 85\,\% & 73 à 94\,\% & 8 à 29\,\% \\
Faux \texttt{conforme} sur les lacunes & 0/13 & 0/13 & 0/13 \\
Intervalle & 0 à 23\,\% & 0 à 23\,\% & 0 à 23\,\% \\
Rappel des lacunes & 13/13 & 13/13 & 13/13 \\
Validité des citations retenues & 100\,\% & 100\,\% & --- \\
Verdicts ramenés à \texttt{non\_trouve} & 2 & 1 & --- \\
Échecs techniques du juge & 0 & 0 & --- \\
\bottomrule
\end{tabular}
\end{center}

Les deux intervalles d'exactitude se recouvrent largement : l'écart de
cinq unités entre les modèles n'est pas significatif sur 44 unités. Les deux
modèles rendent le même verdict final sur 35 unités sur 44.

\subsection*{Par classe de verdict (précision / rappel)}

\begin{center}
\begin{tabular}{@{}lccc@{}}
\toprule
Classe attendue & Effectif & \texttt{qwen3.8:27b} & \texttt{gemma4:31b} \\
\midrule
\texttt{conforme} & 23 & 1,00 / 0,78 & 1,00 / 0,83 \\
\texttt{non\_conforme} & 6 & 0,40 / 1,00 & 0,55 / 1,00 \\
\texttt{non\_trouve} & 7 & 0,60 / 0,43 & 0,83 / 0,71 \\
\texttt{non\_applicable} & 8 & 1,00 / 0,75 & 1,00 / 1,00 \\
\bottomrule
\end{tabular}
\end{center}

\subsection*{Désaccords avec la vérité terrain}

\begin{center}
\begin{tabular}{@{}llll@{}}
\toprule
Unité (\texttt{anx\_II/\ldots}) & Attendu & \texttt{qwen3.8:27b} & \texttt{gemma4:31b} \\
\midrule
\texttt{4/d} & \texttt{non\_trouve} & \texttt{non\_conforme} & \texttt{non\_conforme} \\
\texttt{5/a} & \texttt{non\_trouve} & \texttt{non\_conforme} & \texttt{non\_conforme} \\
\texttt{6.1/a} & \texttt{conforme} & \texttt{non\_conforme} & \texttt{non\_conforme} \\
\texttt{6.1/b/t2} & \texttt{conforme} & \texttt{non\_conforme} & \texttt{non\_conforme} \\
\midrule
\texttt{1.1/g} & \texttt{non\_trouve} & \texttt{non\_conforme} & juste \\
\texttt{1.1/k} & \texttt{conforme} & \texttt{non\_conforme} & juste \\
\texttt{1.2/b} & \texttt{non\_trouve} & \texttt{non\_conforme} & juste \\
\texttt{3/b} & \texttt{conforme} & \texttt{non\_conforme} & juste \\
\texttt{6.1/b/t1} & \texttt{conforme} & \texttt{non\_conforme} & juste \\
\texttt{6.2/c/t2} & \texttt{non\_applicable} & \texttt{non\_trouve}$^{*}$ & juste \\
\texttt{6.2/c/t4} & \texttt{non\_applicable} & \texttt{non\_trouve}$^{*}$ & juste \\
\texttt{1.1/l} & \texttt{conforme} & juste & \texttt{non\_conforme} \\
\texttt{6.1/b/t6} & \texttt{conforme} & juste & \texttt{non\_trouve}$^{*}$ \\
\bottomrule
\end{tabular}
\end{center}

\noindent$^{*}$ Verdict ramené à \texttt{non\_trouve} par la règle de
citation. Pour \texttt{qwen3.8:27b}, le verdict du modèle était
\texttt{non\_applicable} (celui attendu), mais le passage du règlement cité
n'a pas été retrouvé dans l'unité. Pour \texttt{gemma4:31b}, le verdict du
modèle était \texttt{non\_conforme} et deux passages du dossier n'ont pas
été retrouvés.

\section{Constats}

\begin{enumerate}
\item \textbf{La chaîne fonctionne de bout en bout avec un modèle local.}
  88 appels, aucun échec technique, sortie JSON toujours valide avec le
  mode de raisonnement actif, citations retenues toutes vérifiées.
\item \textbf{Aucun faux \texttt{conforme}} sur les 13 lacunes, pour les
  deux modèles. C'est le critère bloquant du projet. L'intervalle (0 à
  23\,\%) rappelle que 13 lacunes ne suffisent pas à garantir un taux bas.
\item \textbf{Toutes les erreurs vont vers la sévérité.} Aucun modèle n'a
  été plus indulgent que l'annotation. L'erreur dominante est de rendre
  \texttt{non\_conforme} là où l'annotation dit \texttt{conforme} ou
  \texttt{non\_trouve}.
\item \textbf{Quatre unités mettent les deux modèles en désaccord avec
  l'annotation, dans le même sens.} Ce sont les premières à relire : soit
  l'annotation est discutable, soit les consignes distinguent mal
  \texttt{non\_trouve} de \texttt{non\_conforme}.
\item \textbf{La règle de citation a un coût mesurable} : elle a annulé
  deux verdicts justes de \texttt{qwen3.8:27b}. C'est le comportement
  voulu, mais il faudra regarder si la citation du règlement échoue pour
  une raison de fond ou de forme.
\end{enumerate}

\section{Limites}

\begin{itemize}
\item Un seul cas, fictif, rédigé et annoté par une IA, non relu. Un juge
  de la même famille de modèles que l'auteur du cas pourrait être avantagé ;
  les deux juges mesurés ici sont d'autres familles.
\item 44 unités, dont 13 lacunes : intervalles larges, aucune hiérarchie
  fiable entre modèles proches.
\item Une seule série par modèle : la variation entre deux séries
  identiques n'est pas mesurée.
\item Le dossier fictif est propre et bien structuré ; ses lacunes sont
  nettes. Un dossier réel, incomplet ou ambigu, est plus difficile.
\item Quantification en 4 bits et mode de raisonnement par défaut : les
  résultats valent pour ces réglages seulement.
\item Le périmètre couvre la présence et la cohérence des éléments
  attendus, pas la qualité technique de leur contenu.
\item Licences des deux modèles non vérifiées.
\end{itemize}

\section{Coût et durée}

\begin{center}
\begin{tabular}{@{}lcc@{}}
\toprule
 & \texttt{qwen3.8:27b} & \texttt{gemma4:31b} \\
\midrule
Durée de la série (44 appels) & 56 min & 48 min \\
Jetons en entrée & 167\,526 & 167\,673 \\
Jetons en sortie & 55\,137 & 37\,053 \\
\bottomrule
\end{tabular}
\end{center}

Machine louée à 0,69 dollar l'heure. Installation et téléchargement des
modèles : environ 15 minutes ; essai préalable sur 3 unités : 6 minutes.
Du lancement à la fin de la seconde série : un peu plus de 2 heures, soit
environ 1,45 dollar, disque compris. Obtenir la machine a demandé sept
tentatives de création, les six premières refusées faute de capacité.

\section{Suites proposées}

\begin{enumerate}
\item Relecture humaine des 13 unités en désaccord, en commençant par les
  quatre communes aux deux modèles, puis des 44 annotations.
\item Un second cas, moins propre, rédigé ou adapté par un humain.
\item Mesure du juge par API sur le même cas, comme point de comparaison.
\item Deux séries identiques d'un même modèle, pour connaître la variation
  de fond avant de comparer des réglages (raisonnement, masquage).
\end{enumerate}

\section*{Fichiers de cette campagne}

\begin{itemize}
\item \texttt{fictif-therm-t1.<modèle>.jsonl} : les 44 prédictions, avec
  citations, passages rejetés et justification du modèle.
\item \texttt{fictif-therm-t1.<modèle>.eval.json} : sortie du banc.
\item \texttt{modeles.txt} : version d'Ollama et fiche des deux modèles.
\end{itemize}

\end{document}
